% !TeX root = ../Masterarbeit.tex

\documentclass[Masterarbeit.tex]{subfile}

\begin{document}
	
	\chapter{Conclusion} \label{conclusions}
	
	This study aims at examining if neural networks, more specifically a partial convolutional U-net, represent a valid alternative to the EnKF data assimilations of NA SPG OHC. Additionally, the utility of this neural network to better estimate the uncertainty of data assimilations is inspected.\\
	The neural network is trained on data assimilations from the well observed Argo-Era (2004 -- 2020), which are assumed to possess less uncertainty due to the integration of more observations. Then, this trained neural network is tasked to reconstruct masked assimilations as well as observations first from its training period, then from earlier times (1958 -- 2004) and finally from 12 newer monthly time steps. The results from these reconstructions are then examined in regard to their agreement to the assimilation as well as other OHC estimates, more specifically those of the IAP and the EN4 objective analysis. On the one hand, this comparison yields conclusions in regard to the neural networks ability to reproduce and replace the data assimilation of three dimensional ocean variables. On the other hand, it helps to estimate the uncertainty of its training data. As the neural network learns the less uncertain assimilation's spatial patterns from well observed time periods, it is not expected that the neural network reconstruction during the pre-Argo-Era agrees entirely with its training data. Instead, the differences between the indirect reconstruction of masked assimilations and the assimilation itself are taken as indicators for uncertainties when the assimilation's model dynamics are not supported by observational input. The differences between the direct neural network  reconstruction of  EN4 observations and the assimilation underline the assimilation's uncertainty at the grid points of observations and its effect on the overall OHC estimate. \\
	
	The results yielded by this approach lead me to the following conclusion:
	\begin{itemize}
		\item Given the close agreement between the neural network reconstruction and the EnKF assimilation I conclude that \textbf{partial convolutional U-nets are able to reproduce EnKF assimilations during its training period, before and after}. In all cases significant correlations in its overall OHC estimates as well as its spatial patterns are found. Even monthly variations in time steps outside of its training data are reproduced by the neural network when directly reconstructing observations from the Argo-Era onwards. Thus, partial convolutional U-nets provide a valid alternative to enhance and reproduce data assimilation of ocean variables. Additionally, the results demonstrate that the neural network reconstructions are also able to replace the EnKF assimilation in their function as a OHC reanalysis for at least a year outside of its training period. \\
		\item Furthermore, this study also demonstrates that neural network reconstructions can take a role in better estimating the uncertainty of EnKF assimilations of OHC. The neural network is able to transfer spatial patterns from its well observed training period to earlier time spans. Consequently, the differences between assimilation and network reconstruction at times with lower observational coverage point towards uncertainties in the assimilation's model dynamics when they are not aided by observational input. This ability is highlighted by the neural network's detection of inconsistencies in the assimilation's misrepresentation of the North Atlantic current's North-West corner at the start of the pre-Argo-Era. The results of the direct reconstruction of observations highlight inconsistencies between assimilation and observations at the observed grid points. Therefore, in addition to replacing them, \textbf{partial convolutional U-nets are able to detect inconsistencies and estimate uncertainties in data assimilation outside of its ensemble spread, especially in regard to time periods with low observational density}.\\
		\item Finally, an analysis of the computational resources used by both the EnKF assimilations and the neural network shows that \textbf{neural networks have the potential to significantly speed up parts of the assimilation process}. The largest differences occur, when the neural network's ability to reconstruct independent of another training instance is utilized. The generation of an additional ensemble member or the reconstruction of adjusted observational data show differences in node hours of up to a factor 1000. With growing resolution and larger datasets and ensembles, this difference between the computational requirements is bound to grow even further (Brune, pers. comm.). \\
	\end{itemize}
	
	With this study, I show that neural networks have the ability to reproduce, extend and even replace assimilations and other mapping methods. Additionally, they are capable of giving insights into the uncertainty of other assimilation techniques, especially for undersampled regions and variables. Therefore, I overall demonstrate that due to their efficiency, flexibility and ability to learn spatial patterns, partial convolutional U-nets can significantly impact the data assimilation of three-dimensional ocean variables.\\
	
	
	
	
	
	
\end{document}